\chapter{Selected Model Case Studies}\label{ch:cases}

Three of the seven audited models are examined here because each shows a different way in which an audit result can mislead a reader who looks only at the headline. They were chosen for what they teach and not as representatives. Per-model values are in Appendix~\ref{app:results}, and nothing here ranks the models.

\section{A fairness risk on a model built to reduce bias}
\code{unitary/unbiased-toxic-roberta} was selected because its card describes training meant to reduce identity bias. The audit raised a fairness risk on it: true-positive rate 0.784 for the unflagged group and 0.692 for the flagged one, a gap of 0.093 with an interval of [0.020, 0.162]. It is one of two models of the seven where the rule fired.

Two points limit what this says. The gap points the opposite way from that of \code{toxic-bert} (Chapter~\ref{ch:example}): here the model catches fewer toxic comments in the flagged group, and its false-positive rate is also higher there (0.106 against 0.034). The method reports the size of a gap and not its direction, so the two cases look alike in a table and differ in kind. And the group is the annotator-reported identity-attack signal, not the identity-term mentions that bias mitigation of this kind targets, so the audit neither confirms nor contradicts the authors' claim. It shows that on this data and grouping recall differs by about nine points, with a lower bound only just above the tolerance (0.0202 against 0.02): a prompt to look further, which is what ``evidence for review'' means.

The model also shows why ratings are not compared across models. Its card and that of \code{toxic-bert} come from the same publisher, have almost the same length (11\,088 and 11\,095 characters) and the same coverage (three of five sections). The LLM rated the first $3/4/3$ and the second $6/7/7$, using different providers (Table~\ref{tab:appratings}), while the deterministic evidence was the same. This is why the consolidated table of Chapter~\ref{ch:results} reports measurements and findings and not scores.

\section{A model audited against the wrong task}
\code{facebook/roberta-hate-speech-dynabench-r4-target} was selected as a case of label-scope mismatch: it detects hate speech, which is narrower than the toxicity labels of Civil Comments. Its profile looks like that of a poor model. Clean accuracy is 0.783; the severe miss rate is 0.837 (interval [0.789, 0.875]), so it caught only 46 of the 282 severe toxic comments. Harmful recall is 0.153 and the false-positive rate modest (0.059). That is the signature of a model with a narrower idea of ``positive'' than the dataset, not of a noisy one. The audit also raised the only robustness risk of the seven (drop 0.072, interval [0.058, 0.087]) and a fairness risk (gap 0.130, interval [0.075, 0.181]). The card is 570 characters, contains none of the five required sections and gives no licence. All five dimensions were rated and the aggregate, 3.65, is the lowest of the five scored audits.

The lesson is about what the numbers mean. TrustLens measures a model against the task it is pointed at, and a mismatch between training objective and task cannot be told from a defect by the probes. The mismatch was known from the selection reason and not discovered by the audit. A reviewer would have to bring that knowledge. A reading that stops at the aggregate also loses the useful part: the loaded weight file matches its published hash, so Integrity is not the problem, whereas the low scores come from behaviour and documentation.

\section{Where the risk rule and the rating part ways}
\code{textdetox/xlmr-large-toxicity-classifier} is a multilingual model listed as an XLM-R large. Three things stand out.

\emph{The risk rule and the rating disagree.} The accuracy drop under typos is 0.054 with an interval of [0.043, 0.065]. The lower bound is under 0.05, so the risk rule raised nothing. The range-scaled mapping, however, gives its lowest rating to any drop above about 0.03, and Robustness is rated $1/1/8$. The two are consistent once the thresholds are known: the rule asks whether the drop is surely above 0.05, the mapping where the drop lies among the defects we injected.

\emph{The fairness rule has a blind spot.} The TPR gap is small (0.074, interval [0.009, 0.138]), so the outcome is ``disparity observed'' and not a risk. The false-positive rates differ far more: 0.167 in the unflagged group and 0.409 in the flagged group, a gap of 0.242, so the model flags 41\% of the benign comments of the flagged group. TrustLens's primary fairness quantity is the TPR gap, so this disparity was recorded as evidence (the equalized-odds difference) but never reached a decision; \code{bert-tiny-toxicity} shows the same pattern (0.350 and 0.553). For a moderation model the over-flagging of one group is a real harm. We did not change the rule after seeing this, since that would be a post-hoc change, and record it as a limitation (Chapter~\ref{ch:discussion}).

\emph{The name does not match the weights.} The repository is called ``large'' and its card describes an XLM-R large, but its weight file has the same hash as that of the sibling repository \code{xlmr-base-toxicity-classifier} and its configuration is of base size. The file check verified it, because the file matches what the Hub publishes for this repository. It cannot notice that the name is misleading, and the discovery came from comparing hashes across repositories by hand. An integrity check says the bytes are the bytes the Hub holds, not that the Hub's description of them is accurate.
